% GENERATED FILE. Edit canonical lessons, then run npm run generate:books.
% canonical-source-hash: fnv1a64:cec988275f511324
% canonical-lessons: JA-C323-enryo-suru, JA-C323-kobamu, JA-C323-muri, JA-C323-kekkou, JA-C323-betsuni

\chapter{Saying No Politely}
\label{ch:ja-saying-no-politely}
\hlchaptermodality{323}

\begin{chapteropening}
\textbf{By the end of this chapter:} \emph{I can turn something down politely and give a reason.}

Complete the last lesson of chapter 323: I can turn something down politely and give a reason.
\end{chapteropening}

\section[\texorpdfstring{\ja{えんりょする} (enryo suru)}{えんりょする (enryo suru)}]{\ja{えんりょする} (enryo suru) — to hold back; to decline politely}

\label{lesson:JA-C323-enryo-suru}



\begin{quote}
Before the new one: say the Japanese for by all means, then the Japanese for to come forward.
\end{quote}

\subsection*{You'll want to know: \ja{えんりょする}}
\textbf{\ja{えんりょする}} — \emph{enryo suru} — \textquotedblleft{}to hold back; to decline politely\textquotedblright{}.

\begin{quote}
\textbf{\ja{あたまが} \ja{いたいので、おさけは} \ja{えんりょします。}} — \emph{atama ga itai node, osake wa enryo shimasu} — \textquotedblleft{}My head hurts, so I will say no to the sake.\textquotedblright{}
\end{quote}

Polite \textbf{\ja{えんりょします}}, \emph{enryoshimasu}; a \ja{する} verb. Past (polite): \textbf{\ja{きゃくは} \ja{えんりょしました。}}, \emph{kyaku wa enryo shimashita}.

\begin{cousinweb}
\textbf{\ja{えんりょする}} (\emph{enryo suru}) is built from a Sino-Japanese noun and the native verb \emph{suru}, \textquotedblleft{}to do\textquotedblright{}. The noun, \emph{enryo}, is two Chinese characters read with their on'yomi, readings borrowed from Chinese: \emph{en}, \textquotedblleft{}far\textquotedblright{} (Middle Chinese \emph{hjwonX}) and \emph{ryo}, \textquotedblleft{}thought\textquotedblright{} (\emph{ljoH}) (KANJIDIC2; JMdict; Middle Chinese in Baxter's transcription).
\end{cousinweb}

\begin{grammarlens}[title={What you've built}]
One more word for giving reasons and opinions.
\end{grammarlens}

\subsection*{Your turn}
\begin{itemize}
\raggedright
  \item \emph{Say it:} \emph{enryo suru}
  \item \emph{Say it:} \emph{enryo suru}, once more
  \item \emph{Say it:} \emph{atama ga itai node, osake wa enryo shimasu} — My head hurts, so I will say no to the sake.
  \item \emph{Recall:} say the Japanese for by all means, then the Japanese for to come forward, then say \emph{enryo suru} again
\end{itemize}

\begin{tcolorbox}[breakable,colback=teal!4,colframe=teal!35!black,title={Before you move on}]
What does \ja{えんりょする} mean? (\textquotedblleft{}To hold back; to decline politely\textquotedblright{}.) Say it once more.
\end{tcolorbox}
\section[\texorpdfstring{\ja{こばむ} (kobamu)}{こばむ (kobamu)}]{\ja{こばむ} (kobamu) — to reject; to refuse firmly}

\label{lesson:JA-C323-kobamu}



\begin{quote}
Before the new one: say the Japanese for to come forward, then the Japanese for to hold back.
\end{quote}

\subsection*{You'll want to know: \ja{こばむ}}
\textbf{\ja{こばむ}} — \emph{kobamu} — \textquotedblleft{}to reject; to refuse firmly\textquotedblright{}.

\begin{quote}
\textbf{\ja{あぶないと} \ja{おもったので、かれは} \ja{しごとを} \ja{こばみました。}} — \emph{abunai to omotta node, kare wa shigoto o kobamimashita} — \textquotedblleft{}He thought it was dangerous, so he refused the job.\textquotedblright{}
\end{quote}

Polite \textbf{\ja{こばみます}}, \emph{kobamimasu}; a godan verb. Past (polite): \textbf{\ja{かれは} \ja{こばみました。}}, \emph{kare wa kobamimashita}.

\begin{cousinweb}
\textbf{\ja{こばむ}} (\emph{kobamu}) is a native Japanese word (wago), inherited rather than borrowed from Chinese, and its older origin is not settled. Its kanji was attached to it later as a kun'yomi, a native reading given to a Chinese character (JMdict, KANJIDIC2).
\end{cousinweb}

\begin{grammarlens}[title={What you've built}]
One more word for giving reasons and opinions.
\end{grammarlens}

\subsection*{Your turn}
\begin{itemize}
\raggedright
  \item \emph{Say it:} \emph{kobamu}
  \item \emph{Say it:} \emph{kobamu}, once more
  \item \emph{Say it:} \emph{abunai to omotta node, kare wa shigoto o kobamimashita} — He thought it was dangerous, so he refused the job.
  \item \emph{Recall:} say the Japanese for to come forward, then the Japanese for to hold back, then say \emph{kobamu} again
\end{itemize}

\begin{tcolorbox}[breakable,colback=teal!4,colframe=teal!35!black,title={Before you move on}]
What does \ja{こばむ} mean? (\textquotedblleft{}To reject; to refuse firmly\textquotedblright{}.) Say it once more.
\end{tcolorbox}
\section[\texorpdfstring{\ja{むり} (muri)}{むり (muri)}]{\ja{むり} (muri) — impossible; unreasonable}

\label{lesson:JA-C323-muri}



\begin{quote}
Before the new one: say the Japanese for to hold back, then the Japanese for to reject.
\end{quote}

\subsection*{You'll want to know: \ja{むり}}
\textbf{\ja{むり}} — \emph{muri} — \textquotedblleft{}impossible; unreasonable\textquotedblright{}.

\begin{quote}
\textbf{\ja{じかんが} \ja{ないから、むりです。}} — \emph{jikan ga nai kara, muri desu} — \textquotedblleft{}There is no time, so it is impossible.\textquotedblright{}
\end{quote}

A na-adjective: \ja{な} before a noun.

\begin{cousinweb}
\textbf{\ja{むり}} (\emph{muri}) is Sino-Japanese (kango): two Chinese characters read with their on'yomi, readings borrowed from Chinese. They are \emph{mu}, \textquotedblleft{}nothing\textquotedblright{} (Middle Chinese \emph{mju}) and \emph{ri}, \textquotedblleft{}reason\textquotedblright{} (\emph{liX}) (KANJIDIC2; Middle Chinese in Baxter's transcription).
\end{cousinweb}

\begin{grammarlens}[title={What you've built}]
One more word for giving reasons and opinions.
\end{grammarlens}

\subsection*{Your turn}
\begin{itemize}
\raggedright
  \item \emph{Say it:} \emph{muri}
  \item \emph{Say it:} \emph{muri}, once more
  \item \emph{Say it:} \emph{jikan ga nai kara, muri desu} — There is no time, so it is impossible.
  \item \emph{Recall:} say the Japanese for to hold back, then the Japanese for to reject, then say \emph{muri} again
\end{itemize}

\begin{tcolorbox}[breakable,colback=teal!4,colframe=teal!35!black,title={Before you move on}]
What does \ja{むり} mean? (\textquotedblleft{}Impossible; unreasonable\textquotedblright{}.) Say it once more.
\end{tcolorbox}
\section[\texorpdfstring{\ja{けっこう} (kekkō)}{けっこう (kekkō)}]{\ja{けっこう} (kekkō) — no thank you; fine}

\label{lesson:JA-C323-kekkou}



\begin{quote}
Before the new one: say the Japanese for to reject, then the Japanese for impossible.
\end{quote}

\subsection*{You'll want to know: \ja{けっこう}}
\textbf{\ja{けっこう}} — \emph{kekkō} — \textquotedblleft{}no thank you; fine\textquotedblright{}.

\begin{quote}
\textbf{\ja{さっき} \ja{たべたから、けっこうです。}} — \emph{sakki tabeta kara, kekkō desu} — \textquotedblleft{}I ate a moment ago, so no, thank you.\textquotedblright{}
\end{quote}

A na-adjective: \ja{な} before a noun.

\begin{cousinweb}
\textbf{\ja{けっこう}} (\emph{kekkō}) is Sino-Japanese (kango): two Chinese characters read with their on'yomi, readings borrowed from Chinese. They are \emph{ketsu}, \textquotedblleft{}tie\textquotedblright{} (Middle Chinese \emph{ket}) and \emph{kō}, \textquotedblleft{}build\textquotedblright{} (\emph{kuwH}) (KANJIDIC2; Middle Chinese in Baxter's transcription). The two join as \emph{kek-kō}.
\end{cousinweb}

\begin{grammarlens}[title={What you've built}]
One more word for giving reasons and opinions.
\end{grammarlens}

\subsection*{Your turn}
\begin{itemize}
\raggedright
  \item \emph{Say it:} \emph{kekkō}
  \item \emph{Say it:} \emph{kekkō}, once more
  \item \emph{Say it:} \emph{sakki tabeta kara, kekkō desu} — I ate a moment ago, so no, thank you.
  \item \emph{Recall:} say the Japanese for to reject, then the Japanese for impossible, then say \emph{kekkō} again
\end{itemize}

\begin{tcolorbox}[breakable,colback=teal!4,colframe=teal!35!black,title={Before you move on}]
What does \ja{けっこう} mean? (\textquotedblleft{}No thank you; fine\textquotedblright{}.) Say it once more.
\end{tcolorbox}
\section[\texorpdfstring{\ja{べつに} (betsuni)}{べつに (betsuni)}]{\ja{べつに} (betsuni) — not particularly}

\label{lesson:JA-C323-betsuni}



\begin{quote}
Before the new one: say the Japanese for impossible, then the Japanese for no thank you.
\end{quote}

\subsection*{You'll want to know: \ja{べつに}}
\textbf{\ja{べつに}} — \emph{betsuni} — \textquotedblleft{}not particularly\textquotedblright{}.

\begin{quote}
\textbf{\ja{べつに} \ja{いそがないから、ゆっくり} \ja{して} \ja{ください。}} — \emph{betsuni isoganai kara, yukkuri shite kudasai} — \textquotedblleft{}There is no particular hurry, so take your time.\textquotedblright{}
\end{quote}

\begin{cousinweb}
\textbf{\ja{べつに}} (\emph{betsuni}) is Sino-Japanese (kango): one Chinese character read with its on'yomi, reading borrowed from Chinese. It is \emph{betsu}, \textquotedblleft{}separate\textquotedblright{} (Middle Chinese \emph{bjet}) (KANJIDIC2; Middle Chinese in Baxter's transcription). The native particle \emph{ni} after it makes the adverb (JMdict).
\end{cousinweb}

\begin{grammarlens}[title={What you've built}]
One more word for giving reasons and opinions.
\end{grammarlens}

\subsection*{Your turn}
\begin{itemize}
\raggedright
  \item \emph{Say it:} \emph{betsuni}
  \item \emph{Say it:} \emph{betsuni}, once more
  \item \emph{Say it:} \emph{betsuni isoganai kara, yukkuri shite kudasai} — There is no particular hurry, so take your time.
  \item \emph{Recall:} say the Japanese for impossible, then the Japanese for no thank you, then say \emph{betsuni} again
\end{itemize}

\begin{tcolorbox}[breakable,colback=teal!4,colframe=teal!35!black,title={Before you move on}]
What does \ja{べつに} mean? (\textquotedblleft{}Not particularly\textquotedblright{}.) Say it once more.
\end{tcolorbox}
